%% kapa: exam=1 round=32 year=2021 period=1 form=A subject=civil_law subject_ko=민법 seq=35 number=35
%% source: Q-Net 감정평가사 기출문제 (cst00309) artlSeq=5208397 file=extracted/1008/5208397/A형/2021년도 제32회 감정평가사 1차 1교시 A형.pdf
%% conversion: kordoc 4.24.51 -> tools/gen1.py (auto); review=proofread-sonnet; flags=none
\begin{question}{35}
2020년 5월 신탁자 甲과 그의 친구인 수탁자 乙이 X부동산에 대하여 명의신탁약정을 한 후, 乙이 직접 계약당사자가 되어 丙으로부터 X를 매수하고 소유권이전등기를 마쳤다. 다음 설명으로 옳지 않은 것은? (다툼이 있으면 판례에 따름)
\begin{choices}
  \item 甲과 乙 사이의 명의신탁약정은 무효이다.
  \item 丙이 甲·乙 사이의 명의신탁약정 사실을 몰랐다면 乙은 X의 소유권을 취득한다.
  \item 丙이 甲·乙 사이의 명의신탁약정 사실을 알았는지 여부는 소유권이전등기가 마쳐진 때를 기준으로 판단하여야 한다.
  \item 乙이 X의 소유자가 된 경우 甲으로부터 제공받은 매수자금 상당액을 甲에게 부당이득으로 반환하여야 한다.
  \item 丙이 甲·乙 사이의 명의신탁약정 사실을 안 경우에도 乙이 그 사정을 모르는 丁에게 X를 매도하여 소유권이전등기를 마쳤다면 丁은 X의 소유권을 취득한다.
\end{choices}
\end{question}
\answer{3}
